\documentclass[10pt,a4paper]{article}
\usepackage[utf8]{inputenc}
\usepackage[english]{babel}
\usepackage{amsmath}
\usepackage{amsfonts}
\usepackage{amssymb}
\usepackage{graphicx}
\usepackage{lmodern}
\usepackage[left=2cm,right=2cm,top=2cm,bottom=2cm]{geometry}
\author{Keith Reid}
\begin{document}
\title{Proposing a Taxonomy of Restriction, Reduction and related Research from a grounded theory involving Needs, Alliance, Shared information, and Help}
\author{Dr Keith S. Reid MBChB MRCPsych LLM PhD}
\maketitle
\section{Abstract}

Meehl and Golden (1982) defined a taxonomy as the construction of classes and the assignment of things into them. 
Reducing Restrictive Interventions comes with many related ideas such as restraint, treatment, advocacy, rights-affirmation, risk assessment, initiatives/policies to influence these, and research.
These may benefit from a taxonomy and the present research group has seen stimulating initial discussion of taxonomies.

Some restraint classifications are valid but trivially dichotomous such as long/short or necessary/unnecessary, and their dividing lines are arguable.
Classifications exist for restraint: ``mechanical", ``cultural", ``chemical"; and for policies, such as the Six Core Strategies.
One risk in taxonomy is restraint-creep, defined here as a social action, suffixing increasingly arbitrary ideas with ``-restraint".
Taking random good things from a list, one may construct culinary restraint, and aesthetic restraint, which sound real but simply bring bland hospital food and decor into a diluted class of ``restraint".
This problem is called a lack of closed-ness under extensibility and reduces how exhaustive, and credible, taxonomies are.

This submission offers a taxonomy, NASH, which can exhaustively classify the concepts in the penumbra while being closed under extensibility.
NASH supposes good care is a co-produced ``negotiation" in a ``dyad" comprising a person in care and their clinician.
The dyad features Needs, Alliance, shared Information, and external Help, hence the acronym NASH.
NASH sits within the present author's theoretical framework, adapanomics, extending the established field of microeconomics.
All relevant restrictions and harms are breakdowns of that negotiation; all policies influence those; all useful research also.
For example, classifying restrictions by the rights they violate is subsumed by Needs or Help in NASH.
NASH allows formulation, interventions, predictions, generalisation, and quantification using every-day language, without creep.

\end{document}